\documentclass[10pt]{amsart}
\usepackage{calc,amssymb,amsmath,ulem,stmaryrd,fullpage}
\usepackage{enumerate}
\usepackage{alltt}
\usepackage{multicol}
\RequirePackage[dvipsnames,usenames]{color}
\let\mffont=\sf
\normalem
%\input{kmacros3.sty}
\input{xy}
\xyoption{all}
\usepackage{hyperref}
\usepackage{graphicx}
\usepackage[all,cmtip]{xy}
\usepackage{verbatim}
\usepackage[left=0.95in,top=0.95in,right=0.95in,bottom=0.95in]{geometry}
\newcommand{\tauCohomology}{T}
\newcommand{\FCohomology}{S}


\theoremstyle{theorem}
%\newtheorem*{mainthm*}{Main Theorem}

\newcommand{\note}[1]{\marginpar{\sffamily\tiny #1}}

\def\todo#1{\textcolor{Mahogany}%
{\sffamily\footnotesize\newline{\color{Mahogany}\fbox{\parbox{\textwidth-15pt}{#1}}}\newline}}

\renewcommand{\O}{\mathcal O}
\newcommand{\ie}{{\itshape i.e.} }
\newcommand{\roundup}[1]{\lceil #1 \rceil}
\newcommand{\rounddown}[1]{\lfloor #1 \rfloor}
\renewcommand{\to}[1][]{\xrightarrow{\ #1\ }}
\newcommand{\onto}[1][]{\protect{\xrightarrow{\ #1\ }\hspace{-0.8em}\rightarrow}}
\newcommand{\into}[1][]{\lhook \joinrel \xrightarrow{\ #1\ }}
%\newcommand{DBDual}[1]{{\underline{\omega}_{#1}}}
\newcommand{\DBDual}[1]{{{\uuline \omega}{}^{\mydot}_{#1}}}
\newcommand{\Tt}{{\mathfrak{T}}}
%\newcommand{\bn}{{\mathfrak{n}} }
\newcommand{\sol}[1]{ \vskip 6pt{\color{blue} {\bf Solution:} #1} \vskip 6pt}

\newcommand{\bR}{{\mathbb{R}}}
\newcommand{\va}{{\vec{a}}}
\newcommand{\vb}{{\vec{b}}}
\newcommand{\vzero}{{\vec{0}}}
\newcommand{\vi}{{\vec{i}}}
\newcommand{\vj}{{\vec{j}}}
\newcommand{\vk}{{\vec{k}}}
\newcommand{\vh}{{\vec{h}}}
\newcommand{\vr}{{\vec{r}}}
\newcommand{\vu}{{\vec{u}}}
\newcommand{\vv}{{\vec{v}}}
\newcommand{\vw}{{\vec{w}}}
\newcommand{\vx}{{\vec{x}}}
\newcommand{\vy}{{\vec{y}}}
\newcommand{\vz}{{\vec{z}}}
\newcommand{\vT}{{\vec{T}}}
\newcommand{\vN}{{\vec{N}}}
\newcommand{\vF}{{\vec{F}}}
\newcommand{\vG}{{\vec{G}}}
\newcommand{\bF}{\mathbb{F}}
\newcommand{\bQ}{\mathbb{Q}}
\newcommand{\bZ}{\mathbb{Z}}
\newcommand{\bC}{\mathbb{C}}


\begin{document}
\title{Homework \#5 -- Math 5320\\Spring 2020}
\author{Due Wednesday, April 15th, in Gradescope}

\maketitle

Note, you can find the book's material on group actions, including solutions to some of these exercises, in section 6.7.  Note I use $.$ to denote the group action, the book uses $*$.  
\noindent
{\bf 1.}  Suppose a group $G$ acts on a set $X$.  Suppose $x \in X$.  Prove that
\[
    \mathrm{stab}_G(x) = \{g \in G\;|\; gx = x\}
\]
is a subgroup of $G$.  
\newpage
\noindent
{\bf 2.}  Suppose a group $G$ acts on a set $X$.  Suppose that $x,y \in X$ and that 
\[
    \mathrm{orbit}_G(x) \bigcap \mathrm{orbit}_G(y) \neq \emptyset.
\]
Prove that $\mathrm{orbit}_G(x) = \mathrm{orbit}_G(y)$.  Then use this show that the distinct orbits partition $X$ (that is, that $X$ is the union of the distint orbits, and distinct orbits have nothing in common).
\vskip 3pt
\emph{Hint:}  Recall that $\mathrm{orbit}_G(x) = \{g.x \;|\; g \in G\} = G.x$
\newpage
\noindent
{\bf 3.}  Suppose a finite group $G$ acts on a set $X$ and that $x \in X$.  Show that 
\[
    |\mathrm{orbit}_G(x)| = |G| / |\mathrm{stab}_G(x)|.
\]   
This result is sometimes called the \emph{orbit-stabilizer theorem}.
\newpage
\noindent
{\bf 4.}  Consider the group $G = S_4$ acting on $\{1, 2, 3, 4\}$.  Show that $H = \mathrm{stab}_G(4)$ is identified with $S_3$.  Now, let $H$ act on $\{1,2,3,4\}$.  What is the orbit of $3$?
\newpage
\noindent
{\bf 5.}  Let $f$ be a polynomial of degree $n$ with coefficients in $F$ and let $K$ be a splitting field for $f$ over $F$.  Prove that $[K : F] \leq n!$.
\newpage
\noindent
{\bf 6.}  Determine the degrees of the splitting fields of the following polynomials over $\bQ$.
\begin{enumerate}
    \item $x^3 - 2$.
    \item $x^6 - 1$.
    \item $x^4 + 1$.
\end{enumerate}
\newpage
\noindent
{\bf 7.}  Let $K = \bQ(\sqrt{2}, \sqrt{3}, \sqrt{5})$.  Compute $[K : \bQ]$.  Prove that $K$ is a Galois extension of $\bQ$.
\newpage
\noindent
{\bf 8.}  Let $K$ be the field from {\bf 7.}  Tell me everything you can about $\mathrm{Aut}(K/\bQ) = G(K/\bQ)$.  For example, is it cyclic? Abelian?  How many elements does it have?
\newpage
\noindent
{\bf 9.}  For each of the following extensions $K \supset \bQ$, determine which are Galois.  If they are not Galois, find a larger extension $K' \supset K$ such that $K' \supset \bQ$ is Galois.
\begin{enumerate}
    \item $\bQ[\sqrt{6}] \supseteq \bQ$.
    \item $\bQ[\sqrt{6}, i] \supseteq \bQ$.
    \item $\bQ[5^{1/4}] \supseteq \bQ$.
\end{enumerate}
\newpage
\noindent
{\bf 10.}  Suppose that $K$ is the splitting field of $x^4 + 5$ over $\bQ$.  Let $H$ be the subgroup of $\mathrm{Aut}(K)$ with two elements, complex conjugation ($a + bi\mapsto a - bi$) and the identity.  Compute the fixed field $K^H$.
\end{document} 